%----------------------------------------------------------------------------------------
% NOTATIE EN FIGUREN VAN DE CURSUS LINEAIRE REGRESSIE
% Staat in de preambule en niet in een hoofdstuk: ook de lijsten met oplossingen gebruiken dit.
% De spreidingsdiagrammen komen uit Spreidingsdiagrammen (2de graad), de dichtheden uit Kansverdelingen;
% nieuw zijn de kwadraten van de residuen en het logaritmisch papier.
%----------------------------------------------------------------------------------------
% Bewaarde grafiek, geplaatst door maak-pdf.py: bestand en pagina in figuren-cache/
\newcommand{\grafiekbestand}[2]{\includegraphics[page=#2,trim=3bp 3bp 3bp 3bp]{figuren-cache/#1}}

\providecommand{\abs}[1]{\lvert #1\rvert}
\providecommand{\R}{\mathbb{R}}
\newcommand{\e}{\mathrm{e}}
\newcommand{\dx}{\,\mathrm{d}x}
% Gemiddelden en voorspelde waarden
\providecommand{\xbar}{\overline{x}}
\providecommand{\ybar}{\overline{y}}
\newcommand{\yhat}{\hat{y}}
% Sommen van kwadraten en producten van de afwijkingen tot het gemiddelde
\newcommand{\Sxx}{S_{xx}}
\newcommand{\Syy}{S_{yy}}
\newcommand{\Sxy}{S_{xy}}
% De totale, verklaarde en onverklaarde kwadratensom (hoofdstuk 2)
\newcommand{\Qt}{Q_{\mathrm{tot}}}
\newcommand{\Qv}{Q_{\mathrm{verkl}}}
\newcommand{\Qf}{Q_{\mathrm{fout}}}
% Partiële afgeleide: \partafg{K}{a} geeft dK/da met ronde d
\newcommand{\partafg}[2]{\frac{\partial #1}{\partial #2}}
\newcommand{\som}{\sum_{i=1}^{n}}
% Compacte matrix in lopende tekst (zoals in Matrices): \km{1 & 2 \\ 3 & 4}
\newcommand{\km}[1]{{\setlength{\arraycolsep}{2.4pt}\renewcommand{\arraystretch}{.82}\left[\begin{array}{*{10}{c}}#1\end{array}\right]}}

% Zoals stappen, maar met een bredere kolom voor de berekening: voor lange formules
\newenvironment{stappenbreed}{%
  \par\renewcommand{\arraystretch}{1.25}%
  \noindent\begin{tabular}{@{}>{$\displaystyle}p{9.8cm}<{$}@{\hspace{4mm}}>{\RaggedRight\small\itshape\arraybackslash}p{4.9cm}@{}}%
}{\end{tabular}\par}

% Zoals stappen, met eigen kolombreedtes: \begin{stappenvrij}{6.2cm}{3.8cm}. Voor naast een figuur.
\newenvironment{stappenvrij}[2]{%
  \par\renewcommand{\arraystretch}{1.25}%
  \noindent\begin{tabular}{@{}>{$\displaystyle}p{#1}<{$}@{\hspace{3mm}}>{\RaggedRight\small\itshape\arraybackslash}p{#2}@{}}%
}{\end{tabular}\par}

% Tekst of berekening links, figuur rechts: \naast{links}{rechts}
\newsavebox{\naastbox}
\newcommand{\naast}[3][]{%
  \par\noindent\sbox{\naastbox}{#3}%
  \begin{minipage}[t]{\dimexpr\linewidth-\wd\naastbox-3mm\relax}\vspace{0pt}#2\end{minipage}\hfill
  \begin{minipage}[t]{\wd\naastbox}\vspace{0pt}\usebox{\naastbox}\end{minipage}\par\addvspace{1mm}}
% Een rij figuren, gecentreerd
\newcommand{\figuren}[1]{\par\addvspace{1mm}{\centering #1\par}\addvspace{1.5mm}}
% Opgave met een figuur rechts: \begin{metgrafiek}{tikz} opgave en deelvragen \end{metgrafiek}
\newsavebox{\grafiekbox}
\newenvironment{metgrafiek}[2][]{%
  \sbox{\grafiekbox}{#2}%
  \begin{minipage}[t]{\dimexpr\linewidth-\wd\grafiekbox-3mm\relax}%
}{%
  \end{minipage}\hfill
  \begin{minipage}[t]{\wd\grafiekbox}\vspace{-2mm}\usebox{\grafiekbox}\end{minipage}\par}

%----------------------------------------------------------------------------------------
% TABELLEN
%----------------------------------------------------------------------------------------
% Tabel met gegevens, horizontaal en met een raster: \waardetabel{aantal kolommen}{x-waarden}{y-waarden},
% de waarden gescheiden door &, in wiskundemodus. Met eigen namen: \waardetabel[$T$ (°C)][$N$]{10}{...}{...}
\NewDocumentCommand{\waardetabel}{O{$x$} O{$y$} m m m}{%
  {\renewcommand{\arraystretch}{1.35}%
  \begin{tabular}{|>{\columncolor{accent!12}}c|*{#3}{>{$}c<{$}|}}\hline #1 & #4 \\ \hline #2 & #5 \\ \hline\end{tabular}}}
% Een kleinere tabel, voor veel kolommen: \kleinetabel[naam x][naam y]{aantal}{x-waarden}{y-waarden}
\NewDocumentCommand{\kleinetabel}{O{$x$} O{$y$} m m m}{%
  {\small\setlength{\tabcolsep}{4pt}\renewcommand{\arraystretch}{1.3}%
  \begin{tabular}{|>{\columncolor{accent!12}}c|*{#3}{>{$}c<{$}|}}\hline #1 & #4 \\ \hline #2 & #5 \\ \hline\end{tabular}}}

%----------------------------------------------------------------------------------------
% GRAFIEKEN MET ASSEN DOOR DE OORSPRONG, binnen een tikzpicture
%----------------------------------------------------------------------------------------
% Rooster met assen: \rooster{xmin}{xstap}{xmax}{ymin}{ystap}{ymax}, met getallen bij elke lijn (niet bij 0).
\NewDocumentCommand{\roosterleeg}{O{$x$} O{$y$} m m m m m m}{%
  \draw[black!14,very thin,xstep=#4,ystep=#7] (#3,#6) grid (#5,#8);
  \draw[->,grijs] (#3,0) -- ({#5+.45*#4},0) node[right,text=black,font=\footnotesize] {#1};
  \draw[->,grijs] (0,#6) -- (0,{#8+.45*#7}) node[above,text=black,font=\footnotesize] {#2};}
\NewDocumentCommand{\rooster}{O{$x$} O{$y$} m m m m m m}{%
  \roosterleeg[#1][#2]{#3}{#4}{#5}{#6}{#7}{#8}
  \pgfmathsetmacro{\rtweex}{#3+#4}\pgfmathsetmacro{\rtweey}{#6+#7}
  \foreach \t in {#3,\rtweex,...,#5}{\pgfmathtruncatemacro{\tt}{round(\t)}\ifdim\t pt=0pt\else\node[below,inner sep=1.5pt,text=black,font=\scriptsize] at (\t,0) {$\tt$};\fi}
  \foreach \t in {#6,\rtweey,...,#8}{\pgfmathtruncatemacro{\tt}{round(\t)}\ifdim\t pt=0pt\else\node[left,inner sep=1.5pt,text=black,font=\scriptsize] at (0,\t) {$\tt$};\fi}}

% Een rechte y = ax + b: \rechte[opties]{xmin}{xmax}{a}{b}
\newcommand{\rechte}[5][]{\draw[accent,thick,#1] (#2,{(#4)*(#2)+(#5)}) -- (#3,{(#4)*(#3)+(#5)});}
% Een grafiek: \kromme[opties]{xmin}{xmax}{voorschrift in \x}
\newcommand{\kromme}[4][]{\draw[accent,thick,#1] plot[domain=#2:#3,samples=90,smooth] (\x,{#4});}
% Een punt: \punt{x}{y}; kleur met de optie
\newcommand{\punt}[3][accent]{\fill[#1] ({#2},{#3}) circle (2.2pt);}
% Een tweede rechte in dezelfde figuur: stijl tweede (gestreept, donker), en hulplijn
\tikzset{tweede/.style={inkt,thick,dashed},hulplijn/.style={grijs,thick,densely dotted}}

%----------------------------------------------------------------------------------------
% SPREIDINGSDIAGRAM, binnen een tikzpicture met x=... en y=... (zoals in Spreidingsdiagrammen)
% De assen beginnen in de hoek (xmin, ymin), niet noodzakelijk in 0.
%----------------------------------------------------------------------------------------
% \diagram[naam x][naam y]{xmin}{xstap}{xmax}{ymin}{ystap}{ymax}: raster en assen met namen.
% Getallen onder en naast de assen: \xas{16,20,...,32} en \yas{0,50,...,250} (komma als decimaalteken).
\NewDocumentCommand{\diagram}{O{$x$} O{$y$} m m m m m m}{%
  \def\dxmin{#3}\def\dymin{#6}%
  \draw[black!14,very thin,xstep=#4,ystep=#7] (#3,#6) grid (#5,#8);
  \draw[->,grijs] (#3,#6) -- ({#5+.5*#4},#6) node[anchor=north east,yshift=-3.6mm,inner sep=1pt,text=black,font=\scriptsize] {#1};
  \draw[->,grijs] (#3,#6) -- (#3,{#8+.5*#7}) node[anchor=south west,xshift=-1mm,inner sep=1pt,text=black,font=\scriptsize] {#2};}
\newcommand{\xas}[1]{\foreach \t in {#1}{\pgfmathtruncatemacro{\hoek}{abs(\t-\dxmin)<0.0001}%
  \ifnum\hoek=1\def\anker{north west}\else\def\anker{north}\fi
  \node[anchor=\anker,inner sep=1.5pt,text=black,font=\scriptsize] at (\t,\dymin) {\pgfmathprintnumber[use comma,1000 sep={}]{\t}};}}
\newcommand{\yas}[1]{\foreach \t in {#1}{\pgfmathtruncatemacro{\hoek}{abs(\t-\dymin)<0.0001}%
  \ifnum\hoek=1\def\anker{south east}\else\def\anker{east}\fi
  \node[anchor=\anker,inner sep=1.5pt,text=black,font=\scriptsize] at (\dxmin,\t) {\pgfmathprintnumber[use comma,1000 sep={}]{\t}};}}
% De gegevens als stippen: \punten{(17,85) (19,88) (20,125)}; een andere kleur met de optie
\newcommand{\punten}[2][accent]{\draw[#1] plot[only marks,mark=*,mark size=1.5pt] coordinates {#2};}
% Het gemiddelde punt (xbar, ybar): een vierkantje in inkt
\newcommand{\gempunt}[2]{\node[fill=inkt,minimum size=4.6pt,inner sep=0pt] at (#1,#2) {};}
% Het residu van een punt tegenover de rechte y = ax + b: \residu{x}{y}{a}{b}, een verticale stippellijn
\newcommand{\residu}[4]{\draw[inkt,thick,densely dashed] (#1,#2) -- (#1,{(#3)*(#1)+(#4)});}
% Het kwadraat van een residu als vierkant naast het residu: \kwadraat{verhouding}{x}{y}{a}{b}.
% De verhouding is (eenheid op de y-as)/(eenheid op de x-as), zodat het vierkant ook op papier vierkant is.
\newcommand{\kwadraat}[5]{\pgfmathsetmacro{\kwr}{(#3)-((#4)*(#2)+(#5))}%
  \fill[fel,fill opacity=.35,draw=fel!70!black,draw opacity=1] (#2,#3) rectangle ({(#2)+(#1)*abs(\kwr)},{(#4)*(#2)+(#5)});}
% De gemiddelden als lijnen door het hele diagram: \gemlijnen{xbar}{ybar}{xmin}{xmax}{ymin}{ymax}
\newcommand{\gemlijnen}[6]{\draw[inkt,thin,dashed] (#1,#5) -- (#1,#6) (#3,#2) -- (#4,#2);}

%----------------------------------------------------------------------------------------
% LOGARITMISCH PAPIER (hoofdstuk 3). In de figuur is de coördinaat de briggse logaritme van de waarde.
%----------------------------------------------------------------------------------------
% Horizontale lijnen van een logaritmische y-as: \logy{xmin}{xmax}{eerste decade}{laatste decade},
% met de decades als exponent (0 is 1, 1 is 10, ...): lijnen bij 1, 2, ..., 9 keer elke macht van 10.
\newcommand{\logy}[4]{%
  \foreach \d in {#3,...,#4}{\foreach \m in {1,...,9}{%
    \pgfmathsetmacro{\hh}{\d+log10(\m)}\pgfmathtruncatemacro{\binnen}{(\hh<=#4+0.0001)?1:0}%
    \ifnum\m=1 \def\logtint{30}\else\def\logtint{12}\fi
    \ifnum\binnen=1 \draw[black!\logtint,very thin] (#1,\hh) -- (#2,\hh);\fi}}}
% Verticale lijnen van een logaritmische x-as: \logx{ymin}{ymax}{eerste decade}{laatste decade}
\newcommand{\logx}[4]{%
  \foreach \d in {#3,...,#4}{\foreach \m in {1,...,9}{%
    \pgfmathsetmacro{\hh}{\d+log10(\m)}\pgfmathtruncatemacro{\binnen}{(\hh<=#4+0.0001)?1:0}%
    \ifnum\m=1 \def\logtint{30}\else\def\logtint{12}\fi
    \ifnum\binnen=1 \draw[black!\logtint,very thin] (\hh,#1) -- (\hh,#2);\fi}}}

%----------------------------------------------------------------------------------------
% DICHTHEDEN (hoofdstuk 2), zoals in Kansverdelingen. In TikZ is \x de veranderlijke:
%   gauss(x, mu, sigma)  normale dichtheid
%   tdich(x, nu)         t-verdeling met nu vrijheidsgraden (nu geheel; exact tot 12, daarboven op 0,05 % na)
%----------------------------------------------------------------------------------------
\tikzset{declare function={
  gauss(\x,\m,\s)=exp(-((\x-\m)^2)/(2*\s^2))/(\s*sqrt(2*pi));
  tcon(\n)=(\n==1)*0.318310+(\n==2)*0.353553+(\n==3)*0.367553+(\n==4)*0.375+(\n==5)*0.379607+(\n==6)*0.382733
    +(\n==7)*0.384991+(\n==8)*0.386699+(\n==9)*0.388035+(\n==10)*0.389108+(\n==11)*0.389990+(\n==12)*0.390726
    +(\n>12)*0.398942*(1-1/(4*\n)+1/(32*\n*\n));
  tdich(\x,\n)=tcon(\n)*exp(-(\n+1)/2*ln(1+\x*\x/\n));
}}
% Assen voor een dichtheid: \kansassen{xmin}{xmax}{ymax}{label x}
\newcommand{\kansassen}[4]{%
  \draw[->,grijs] (#1,0) -- (#2,0) node[right,text=black] {#4};
  \draw[->,grijs] (#1,0) -- (#1,#3);}
